% Source: https://artofproblemsolving.com/community/c3h2590352_trapezium_prob
% Stable identity: aops-topic-2590352-q01
Let $AB=b$, $CD=c$ and let $h$ be the height of the trapezoid. Use coordinates $A=(0,0)$, $B=(b,0)$, $D=(d,h)$, $C=(d+c,h)$. Then
\[P=\left(\dfrac b2,0\right),\quad Q=\left(\dfrac{b+d}2,\dfrac h2\right),\quad R=\left(d+\dfrac c2,h\right),\quad S=\left(\dfrac{d+c}2,\dfrac h2\right).\]
$Q$ and $S$ lie on the midline $y=\dfrac h2$, with $QS=\dfrac{|b-c|}{2}$, while $P$ (on $y=0$) and $R$ (on $y=h$) lie on opposite sides of it, each at distance $\dfrac h2$. Splitting $PQRS$ along $QS$,
\[[PQRS]=\dfrac12\cdot QS\cdot\dfrac h2+\dfrac12\cdot QS\cdot\dfrac h2=\dfrac{|b-c|\,h}{4}.\]

Triangles $AOB$ and $COD$ are similar with ratio $b:c$, so $O$ divides the height in the ratio $b:c$. Its distance to $AB$ is $\dfrac{bh}{b+c}$ and to $CD$ it is $\dfrac{ch}{b+c}$, hence
\[[AOB]=\dfrac{b^2h}{2(b+c)},\qquad[COD]=\dfrac{c^2h}{2(b+c)},\qquad[COD]-[AOB]=\dfrac{\left(c^2-b^2\right)h}{2(b+c)}=\dfrac{(c-b)h}{2}.\]
This difference is $2020-2019=1$, so $c>b$ and $(c-b)h=2$. Therefore
\[[PQRS]=\dfrac{(c-b)h}{4}=\boxed{\dfrac12}.\]

\begin{asy}
usepackage("amsmath");
size(6cm);
defaultpen(fontsize(10pt));
pair A=(0,0),B=(3,0),D=(1,3),C=(8,3);
pair O=extension(A,C,B,D);
pair P=(A+B)/2,Q=(B+D)/2,R=(D+C)/2,Ss=(C+A)/2;
draw(A--B--C--D--cycle,linewidth(0.8));draw(A--C,gray);draw(B--D,gray);
filldraw(P--Q--R--Ss--cycle,lightgray,linewidth(0.8));
draw((-0.3,1.5)--(8.3,1.5),dashed);
dot(O);dot(P);dot(Q);dot(R);dot(Ss);
label("$A$",A,SW);label("$B$",B,SE);label("$C$",C,NE);label("$D$",D,NW);label("$O$",O,S);
label("$P$",P,S);label("$Q$",Q,NW);label("$R$",R,N);label("$S$",Ss,SE);
\end{asy}

The sketch is not to scale (with the given areas, $AB$ and $CD$ are nearly equal and $PQRS$ is very thin). $Q$ and $S$ lie on the dashed midline.
